\documentclass[12pt,a4paper,french]{report}

\usepackage[utf-8]{inputenc}
\usepackage[T1]{fontenc}
\usepackage[french]{babel}
\usepackage{geometry}
\usepackage{fancyhdr}
\usepackage{graphicx}
\usepackage{tikz}
\usepackage{xcolor}
\usepackage{tabularx}
\usepackage{booktabs}
\usepackage{hyperref}
\usepackage{tocloft}
\usepackage{listings}
\usepackage{array}
\usepackage{datetime}
\usepackage{float}

% Configuration de la géométrie
\geometry{
    left=2cm,
    right=2cm,
    top=2.5cm,
    bottom=2.5cm,
    headheight=15pt
}

% Couleurs du projet
\definecolor{SolarBlue}{RGB}{0, 112, 192}
\definecolor{SolarGreen}{RGB}{34, 139, 34}
\definecolor{SolarOrange}{RGB}{255, 140, 0}
\definecolor{DarkGray}{RGB}{64, 64, 64}
\definecolor{LightGray}{RGB}{240, 240, 240}

% Configuration du header et footer
\pagestyle{fancy}
\fancyhf{}
\fancyhead[L]{\small Charte du Projet SolarEase}
\fancyhead[R]{\small Version 1.0}
\fancyfoot[C]{\thepage}
\renewcommand{\headrulewidth}{0.5pt}
\renewcommand{\footrulewidth}{0.5pt}

% Configuration des sections
\usepackage{titlesec}
\titleformat{\chapter}[hang]{\Large\bfseries\color{SolarBlue}}{\thechapter.}{10pt}{}
\titleformat{\section}[hang]{\large\bfseries\color{SolarBlue}}{\thesection.}{10pt}{}
\titleformat{\subsection}[hang]{\normalsize\bfseries\color{DarkGray}}{\thesubsection.}{10pt}{}

% Configuration des listes
\usepackage{enumitem}
\setlist[itemize]{label=\textbullet, leftmargin=2em}
\setlist[enumerate]{leftmargin=2em}

% Configuration du code
\lstset{
    basicstyle=\ttfamily\small,
    breaklines=true,
    breakatwhitespace=true,
    frame=single,
    backgroundcolor=\color{LightGray},
    keywordstyle=\color{SolarBlue}\bfseries,
    commentstyle=\color{gray},
    stringstyle=\color{SolarGreen},
    columns=flexible,
    numbers=left,
    numberstyle=\tiny,
    tabsize=4
}

% Nouvelle commande pour les boîtes de contenu
\newcommand{\boitetitre}[1]{
    \colorbox{SolarBlue}{\begin{minipage}{\textwidth}
    \textcolor{white}{\textbf{#1}}
    \end{minipage}}
}

\newcommand{\boxcontent}[2]{
    \begin{table}[H]
    \centering
    \begin{tabularx}{\textwidth}{|X|}
    \hline
    \rowcolor{LightGray}
    \textbf{#1} \\
    \hline
    #2 \\
    \hline
    \end{tabularx}
    \end{table}
}

%================= DÉBUT DU DOCUMENT =================

\begin{document}

% =========== PAGE DE TITRE ===========
\begin{titlepage}
    \centering
    \vspace*{2cm}
    
    % Logo ou titre principal
    \begin{tikzpicture}
        \draw[fill=SolarBlue,draw=SolarGreen, line width=2pt] (0,0) rectangle (6,2);
        \node at (3,1) {\huge\textcolor{white}{\textbf{SOLAREASE}}};
    \end{tikzpicture}
    
    \vspace{1.5cm}
    
    {\Large \textbf{\textcolor{SolarBlue}{CHARTE DE PROJET}}}
    
    \vspace{2cm}
    
    {\LARGE \textbf{Plateforme de Gestion Intelligente}}
    
    \vspace{0.5cm}
    
    {\LARGE \textbf{de l'Énergie Solaire}}
    
    \vspace{3cm}
    
    \begin{tabularx}{\textwidth}{|X|X|}
    \hline
    \rowcolor{LightGray}
    \textbf{Version} & 1.0 \\
    \hline
    \textbf{Date} & \today \\
    \hline
    \textbf{Statut} & Approuvé \\
    \hline
    \textbf{Classification} & Confidentiel \\
    \hline
    \end{tabularx}
    
    \vspace{3cm}
    
    \begin{tabularx}{\textwidth}{|l|X|}
    \hline
    \rowcolor{LightGray}
    \multicolumn{2}{|c|}{\textbf{INFORMATIONS DE CONTACT}} \\
    \hline
    \textbf{Chef de Projet} & À compléter \\
    \hline
    \textbf{Sponsor} & À compléter \\
    \hline
    \textbf{Responsable Technique} & À compléter \\
    \hline
    \end{tabularx}
    
    \vfill
    
    {\small Année Académique : 2025-2026}
\end{titlepage}

% =========== TABLE DES MATIÈRES ===========
\cleardoublepage
\tableofcontents
\cleardoublepage

% =========== CHAPITRE 1: INTRODUCTION ===========
\chapter{Introduction}

\section{Contexte du Projet}

SolarEase est une plateforme intégrée de gestion intelligente de l'énergie solaire, conçue pour optimiser l'utilisation des installations photovoltaïques. Le projet répond à une demande croissante de solutions numériques dans le secteur des énergies renouvelables, en particulier pour la gestion, le dimensionnement et le suivi des projets solaires.

\subsection{Présentation du Projet PFE}

Mon projet principal est SolarEase, réalisé dans le cadre de mon PFE. C'est une plateforme de dimensionnement solaire construite en architecture microservices avec Spring Boot côté backend et React/TypeScript côté frontend. J'y ai travaillé sur l'authentification, les APIs REST, la logique métier de dimensionnement, et l'expérience utilisateur entre espace client et espace installateur. J'ai aussi contribué à l'intégration de l'IA avec une approche RAG, pour enrichir les recommandations techniques et financières à partir d'une base de connaissances. L'objectif était d'aider à la décision avec des réponses plus pertinentes et contextualisées.

\textbf{Déploiement en production} : Le frontend a été déployé avec \textbf{Vercel}, offrant une distribution globale et des performances optimales. Le backend et la base de données ont été déployés sur \textbf{Railway}, assurant une infrastructure cloud stable et scalable pour les services microservices.

\subsection{Enjeux}

\begin{itemize}
    \item \textbf{Énergétique} : Optimiser la production et la consommation d'énergie solaire
    \item \textbf{Économique} : Réduire les coûts opérationnels et améliorer le ROI
    \item \textbf{Environnemental} : Promouvoir l'adoption des énergies renouvelables
    \item \textbf{Technologique} : Fournir une solution moderne et scalable
\end{itemize}

\section{Vision et Objectifs du Projet}

\subsection{Vision}

Devenir la plateforme de référence pour la gestion intelligente et intégrée des installations solaires, offrant des outils puissants et intuitifs aux professionnels et aux gestionnaires d'énergie.

\subsection{Objectifs Globaux}

\begin{enumerate}
    \item Créer une application web moderne et ergonomique
    \item Développer des services microservices robustes et scalables
    \item Implémenter des fonctionnalités avancées de dimensionnement solaire
    \item Assurer la sécurité et l'authenticité des utilisateurs
    \item Permettre la gestion complète du cycle de vie des projets solaires
\end{enumerate}

% =========== CHAPITRE 2: DÉFINITION DU PROJET ===========
\chapter{Définition du Projet}

\section{Périmètre du Projet}

\subsection{Inclusions}

\boitetitre{Composants Inclus}

\begin{itemize}
    \item \textbf{Frontend} : Application web React avec interface utilisateur moderne
    \item \textbf{Backend} : Architecture microservices (4 services + Gateway)
    \item \textbf{Authentification} : Système d'authentification JWT sécurisé
    \item \textbf{Base de données} : PostgreSQL pour la persistance des données
    \item \textbf{Infrastructure} : Dockerisation et déploiement containerisé
    \item \textbf{Documentation} : API REST, guides utilisateurs, documentation technique
\end{itemize}

\subsection{Exclusions}

\begin{itemize}
    \item Installation physique des panneaux solaires
    \item Support hardware des capteurs IoT en phase 1
    \item Intégration avec des systèmes tiers (calendrier, facturations externes, etc.)
    \item Formation avancée pour utilisateurs externes (phase 2)
\end{itemize}

\section{Description des Composantes}

\subsection{Architecture Technique}

SolarEase est basée sur une architecture microservices moderne avec les composantes suivantes :

\subsubsection{1. Frontend - React Application}

\begin{itemize}
    \item \textbf{Framework} : React 18+ avec Vite
    \item \textbf{Design System} : Material-UI (MUI) + Radix UI
    \item \textbf{Styling} : Tailwind CSS, PostCSS
    \item \textbf{Gestion d'État} : Context API avec React Hooks
    \item \textbf{Type Safety} : TypeScript
    \item \textbf{Build Tool} : Vite pour performances optimales
\end{itemize}

\subsubsection{2. Services Microservices (Java Spring Boot)}

\paragraph{a) Gateway Service}
\begin{itemize}
    \item Point d'entrée unique pour tous les appels API
    \item Gestion du routage vers les différents services
    \item Contrôle des autorisations et middleware
    \item Version : 1.0.0
\end{itemize}

\paragraph{b) Identity Service}
\begin{itemize}
    \item Gestion de l'authentification et autorisation
    \item Gestion des utilisateurs et des rôles
    \item Génération et validation des tokens JWT
    \item Version : 1.0.0
\end{itemize}

\paragraph{c) Dimensioning Service}
\begin{itemize}
    \item Calculs avancés de dimensionnement solaire
    \item Analyse des toitures et orientations
    \item Estimations de production énergétique
    \item Optimisation des configurations
    \item Version : 1.0.0
\end{itemize}

\paragraph{d) Project Service}
\begin{itemize}
    \item Gestion du cycle de vie des projets
    \item Gestion des équipements et installations
    \item Suivi des clients et contrats
    \item Gestion des statuts et historique
    \item Version : 1.0.0
\end{itemize}

\subsubsection{3. Infrastructure}

\begin{itemize}
    \item \textbf{Containerization} : Docker et Docker Compose
    \item \textbf{Base de données} : PostgreSQL
    \item \textbf{Runtime} : Java 17 (services backend)
    \item \textbf{Node.js} : v18+ (frontend)
\end{itemize}

\section{Livrables Principaux}

\begin{table}[H]
\centering
\resizebox{\textwidth}{!}{
\begin{tabularx}{\textwidth}{|l|l|X|l|}
\hline
\rowcolor{LightGray}
\textbf{N°} & \textbf{Livrable} & \textbf{Description} & \textbf{Statut} \\
\hline
1 & Application Frontend & Interface web complète SolarEase & En cours \\
\hline
2 & Gateway Service & Service de routage API & Complété \\
\hline
3 & Identity Service & Service d'authentification & Complété \\
\hline
4 & Dimensioning Service & Service de calcul solaire & Complété \\
\hline
5 & Project Service & Service de gestion de projets & Complété \\
\hline
6 & Base de données & Schéma PostgreSQL complet & Complété \\
\hline
7 & Documentation API & Spécifications OpenAPI/Swagger & En cours \\
\hline
8 & Documentation Technique & Guides d'installation et déploiement & Complété \\
\hline
9 & Tests d'Intégration & Suite complète de tests unitaires & Complété \\
\hline
10 & Guide Utilisateur & Documentation fonctionnelle & Planifié \\
\hline
\end{tabularx}
}
\end{table}

% =========== CHAPITRE 3: ORGANISATION ===========
\chapter{Organisation du Projet}

\section{Structure Organisationnelle}

\subsection{Rôles et Responsabilités}

\begin{table}[H]
\centering
\begin{tabularx}{\textwidth}{|l|X|X|}
\hline
\rowcolor{LightGray}
\textbf{Rôle} & \textbf{Responsabilités} & \textbf{Titulaire} \\
\hline
\textbf{Chef de Projet} & Pilotage général, gestion des risques, coordination & À désigner \\
\hline
\textbf{Sponsor Exécutif} & Approbation des décisions majeures, budget & À désigner \\
\hline
\textbf{Architecte Technique} & Architecture, décisions techniques majeures & Équipe Dev \\
\hline
\textbf{Lead Developer Frontend} & Supervision frontend, qualité code UI/UX & Équipe Frontend \\
\hline
\textbf{Lead Developer Backend} & Supervision services, architecture microservices & Équipe Backend \\
\hline
\textbf{Responsable QA} & Tests, qualité, documentation tests & Équipe QA \\
\hline
\textbf{DevOps Engineer} & Déploiement, infrastructure, CI/CD & Équipe Ops \\
\hline
\end{tabularx}
\end{table}

\section{Structure du Projet}

\subsection{Packages et Dépendances}

Le projet est organisé selon cette structure :

\begin{lstlisting}[language=bash]
projetPfe/
├── Design Authentication Flow/         # Prototype authentification
│   └── src/app/pages/
├── SolarEase Web App Design/           # Application React principale
│   ├── src/
│   │   ├── main.tsx
│   │   ├── app/
│   │   │   ├── components/
│   │   │   ├── pages/
│   │   │   ├── services/
│   │   │   └── context/
│   │   └── styles/
│   └── package.json
├── gateway-service/                     # API Gateway Spring Boot
│   └── pom.xml
├── identity-service/                    # Service d'authentification
│   └── pom.xml
├── dimensioning-service/                # Service de dimensionnement
│   └── pom.xml
├── project-service/                     # Service de gestion projets
│   └── pom.xml
└── init.sql/                            # Initialisation DB PostgreSQL
\end{lstlisting}

% =========== CHAPITRE 4: PLANNING ===========
\chapter{Planning et Échéancier}

\section{Phases du Projet}

\subsection{Phase 1 : Préparation et Setup (Semaines 1-2)}

\begin{itemize}
    \item Mise en place de l'environnement de développement
    \item Configuration des outils et infrastructure
    \item Réunion de lancement et alignement
    \item Définition des conventions de code
\end{itemize}

\subsection{Phase 2 : Développement Backend (Semaines 2-6)}

\begin{itemize}
    \item Implémentation de Gateway Service
    \item Implémentation de Identity Service
    \item Implémentation de Dimensioning Service
    \item Implémentation de Project Service
    \item Tests unitaires et d'intégration
\end{itemize}

\subsection{Phase 3 : Développement Frontend (Semaines 4-8)}

\begin{itemize}
    \item Implémentation des pages d'authentification
    \item Développement de l'interface utilisateur
    \item Intégration avec les APIs backend
    \item Tests de l'interface et UX
\end{itemize}

\subsection{Phase 4 : Intégration et Tests (Semaines 7-9)}

\begin{itemize}
    \item Tests d'intégration système complets
    \item Tests de performance et charge
    \item Corrective des bugs identifiés
    \item Tests de sécurité
\end{itemize}

\subsection{Phase 5 : Déploiement et Livraison (Semaines 9-10)}

\begin{itemize}
    \item Mise en production de l'application
    \item Documentation finale
    \item Formation utilisateurs
    \item Support initial
\end{itemize}

\section{Calendrier Détaillé}

\begin{table}[H]
\centering
\resizebox{\textwidth}{!}{
\begin{tabularx}{\textwidth}{|l|l|l|l|}
\hline
\rowcolor{LightGray}
\textbf{Jalons} & \textbf{Dates} & \textbf{Livrables} & \textbf{Statut} \\
\hline
Kick-off & Week 1 & Charter approuvée & ✓ \\
\hline
Setup Environnement & Week 2 & Stack complète opérationnelle & ✓ \\
\hline
Backend v1 & Week 6 & Services testés et documentés & ✓ \\
\hline
Frontend v1 & Week 8 & Pages principales intégrées & En cours \\
\hline
Intégration Complète & Week 9 & Tests d'intégration passés & Planifié \\
\hline
Mise en Production & Week 10 & Application en production & Planifié \\
\hline
\end{tabularx}
}
\end{table}

% =========== CHAPITRE 5: RESSOURCES ===========
\chapter{Ressources et Budget}

\section{Ressources Humaines}

\subsection{Équipe}

\begin{itemize}
    \item \textbf{1 Chef de Projet} : Pilotage et suivi
    \item \textbf{2 Développeurs Frontend} : React, TypeScript, UI/UX
    \item \textbf{2 Développeurs Backend} : Java, Spring Boot, Microservices
    \item \textbf{1 DevOps} : Docker, déploiement, infrastructure
    \item \textbf{1 Testeur QA} : Tests et validation
    \item \textbf{1 Architecte Technique} : Conseil et validation architectural
\end{itemize}

\subsection{Compétences Requises}

\begin{itemize}
    \item \textbf{Frontend} : React, TypeScript, Material-UI, Tailwind CSS
    \item \textbf{Backend} : Java 17, Spring Boot, Spring Cloud, JPA/Hibernate
    \item \textbf{DevOps} : Docker, Docker Compose, PostgreSQL
    \item \textbf{Tests} : JUnit, Mockito, Selenium (si nécessaire)
    \item \textbf{Méthodologies} : Agile, Git/GitHub, CI/CD
\end{itemize}

\section{Ressources Technologiques}

\subsection{Stack Technique}

\begin{table}[H]
\centering
\begin{tabularx}{\textwidth}{|l|l|l|}
\hline
\rowcolor{LightGray}
\textbf{Domaine} & \textbf{Technologie} & \textbf{Version} \\
\hline
\textbf{Frontend} & React & 18.x \\
\hline
 & TypeScript & 5.x \\
\hline
 & Material-UI & 7.3.5 \\
\hline
 & Tailwind CSS & Latest \\
\hline
 & Vite & Latest \\
\hline
\textbf{Backend} & Java & 17 \\
\hline
 & Spring Boot & 3.2.1 \\
\hline
 & Spring Cloud & 2023.0.0 \\
\hline
\textbf{Database} & PostgreSQL & 13+ \\
\hline
\textbf{DevOps} & Docker & 24.x \\
\hline
 & Docker Compose & 2.x \\
\hline
\textbf{Outils} & Maven & 3.8+ \\
\hline
 & npm & 10.x \\
\hline
 & Git & 2.40+ \\
\hline
\end{tabularx}
\end{table}

\section{Budget Estimé}

\begin{table}[H]
\centering
\begin{tabularx}{\textwidth}{|l|X|l|l|}
\hline
\rowcolor{LightGray}
\textbf{Catégorie} & \textbf{Description} & \textbf{Coût Unitaire} & \textbf{Total} \\
\hline
\textbf{Ressources Humaines} & 8 personnes × 10 semaines & Variable & Voir détails \\
\hline
\textbf{Infrastructure} & Serveurs, hosting cloud & À définir & À définir \\
\hline
\textbf{Outils Logiciels} & Licences, services cloud & À définir & À définir \\
\hline
\textbf{Formation} & Formation équipe & À définir & À définir \\
\hline
\textbf{Contingence} & Réserve (15\%) & À définir & À définir \\
\hline
\textbf{TOTAL} & & & À définir \\
\hline
\end{tabularx}
\end{table}

% =========== CHAPITRE 6: RISQUES ===========
\chapter{Gestion des Risques}

\section{Identification des Risques}

\subsection{Risques Techniques}

\begin{table}[H]
\centering
\begin{tabularx}{\textwidth}{|l|l|l|l|}
\hline
\rowcolor{LightGray}
\textbf{Risque} & \textbf{Probabilité} & \textbf{Impact} & \textbf{Mitigation} \\
\hline
Problèmes de performance & Moyenne & Élevé & Optimisation, load testing \\
\hline
Sécurité des données & Moyenne & Critique & Audit sécu, chiffrement \\
\hline
Incompatibilité API & Faible & Moyen & Tests d'intégration précoce \\
\hline
Scalabilité insuffisante & Faible & Élevé & Architecture microservices \\
\hline
\end{tabularx}
\end{table}

\subsection{Risques Organisationnels}

\begin{table}[H]
\centering
\begin{tabularx}{\textwidth}{|l|l|l|l|}
\hline
\rowcolor{LightGray}
\textbf{Risque} & \textbf{Probabilité} & \textbf{Impact} & \textbf{Mitigation} \\
\hline
Manque de ressources & Faible & Élevé & Planification claire, buffer \\
\hline
Retards d'équipe & Moyenne & Moyen & Suivi régulier, Agile \\
\hline
Communication insuffisante & Moyenne & Moyen & Réunions régulières \\
\hline
Turnover d'équipe & Faible & Critique & Engagement, documentation \\
\hline
\end{tabularx}
\end{table}

\section{Stratégies d'Atténuation}

\begin{itemize}
    \item \textbf{Technique} : Code reviews, tests automatisés, documentation
    \item \textbf{Processus} : Sprints courts, retrospectives, amélioration continue
    \item \textbf{Communication} : Daily standups, sprint reviews, stakeholder updates
    \item \textbf{Qualité} : QA rigoureuse, performance tests, security audits
\end{itemize}

% =========== CHAPITRE 7: QUALITÉ ===========
\chapter{Assurance Qualité}

\section{Critères de Qualité}

\subsection{Code}

\begin{itemize}
    \item \textbf{Coverage} : Minimum 80\% de couverture de tests
    \item \textbf{Complexité} : Cyclomatic complexity < 10
    \item \textbf{Duplication} : < 5\% code dupliqué
    \item \textbf{Standards} : Respect des conventions de code
\end{itemize}

\subsection{Performance}

\begin{itemize}
    \item \textbf{Temps de réponse} : < 500ms pour la majorité des requests
    \item \textbf{Disponibilité} : SLA 99.5\%
    \item \textbf{Scalabilité} : Support de 1000+ utilisateurs simultanés
\end{itemize}

\subsection{Sécurité}

\begin{itemize}
    \item \textbf{Authentification} : JWT tokens sécurisés
    \item \textbf{Autorisation} : RBAC (Role-Based Access Control)
    \item \textbf{Encryption} : HTTPS/TLS, données chiffrées en transit et au repos
    \item \textbf{Validation} : Validation des inputs, protection SQL injection
\end{itemize}

\section{Stratégie de Tests}

\subsection{Tests Unitaires}

\begin{itemize}
    \item Chaque composant React testé avec Jest et React Testing Library
    \item Chaque service Java testé avec JUnit et Mockito
    \item Objectif : 80\%+ coverage
\end{itemize}

\subsection{Tests d'Intégration}

\begin{itemize}
    \item Tests des flux complets API
    \item Tests de l'intégration frontend-backend
    \item Tests de la base de données
\end{itemize}

\subsection{Tests de Performance}

\begin{itemize}
    \item Load testing avec JMeter
    \item Profiling des endpoints critiques
    \item Optimisation des requêtes DB
\end{itemize}

\subsection{Tests de Sécurité}

\begin{itemize}
    \item Audit de sécurité du code
    \item Tests de pénétration
    \item Validation des tokens JWT
    \item Tests d'autorisation
\end{itemize}

% =========== CHAPITRE 8: COMMUNICATION ===========
\chapter{Communication et Gouvernance}

\section{Plan de Communication}

\subsection{Réunions Régulières}

\begin{table}[H]
\centering
\begin{tabularx}{\textwidth}{|l|l|l|X|}
\hline
\rowcolor{LightGray}
\textbf{Réunion} & \textbf{Fréquence} & \textbf{Durée} & \textbf{Participants} \\
\hline
Daily Standup & Quotidienne & 15 min & Équipe complète \\
\hline
Sprint Planning & Bi-hebdomadaire & 2h & Équipe + PM \\
\hline
Sprint Review & Bi-hebdomadaire & 1h & Équipe + Stakeholders \\
\hline
Sprint Retrospective & Bi-hebdomadaire & 1h & Équipe \\
\hline
Status Report & Hebdomadaire & 30 min & PM + Sponsor \\
\hline
\end{tabularx}
\end{table}

\section{Escalade des Problèmes}

\begin{enumerate}
    \item \textbf{Niveau 1} : Équipe de développement (résolution dans 24h)
    \item \textbf{Niveau 2} : Chef de Projet + Tech Lead (résolution dans 48h)
    \item \textbf{Niveau 3} : Sponsor + Architecte (résolution immédiate)
\end{enumerate}

\section{Documentation}

\begin{itemize}
    \item \textbf{Wiki du projet} : Confluence ou GitHub Wiki
    \item \textbf{API Documentation} : Swagger/OpenAPI
    \item \textbf{Code Documentation} : Javadoc, JSDoc, README
    \item \textbf{Guides utilisateurs} : Manuels PDF et vidéos
    \item \textbf{Architecture Documentation} : Diagrammes C4, ADRs
\end{itemize}

% =========== CHAPITRE 9: SUCCÈS ===========
\chapter{Critères de Succès}

\section{Objectifs Mesurables}

\subsection{Fonctionnalité}

\begin{itemize}
    \item ✓ Tous les user stories planifiés implémentés
    \item ✓ API conformes aux spécifications OpenAPI
    \item ✓ Authentification fonctionnelle et sécurisée
    \item ✓ Dimensionnement solaire opérationnel
    \item ✓ Gestion de projets complète
\end{itemize}

\subsection{Qualité}

\begin{itemize}
    \item ✓ 80\%+ couverture de tests
    \item ✓ 0 bugs critiques en production
    \item ✓ Performance < 500ms en moyenne
    \item ✓ Disponibilité 99.5\%
\end{itemize}

\subsection{Délais}

\begin{itemize}
    \item ✓ Livraison en production dans les 10 semaines
    \item ✓ Respect du budget alloué
    \item ✓ Aucun dépassement majeur
\end{itemize}

\subsection{Satisfaction}

\begin{itemize}
    \item ✓ Satisfaction utilisateur ≥ 8/10
    \item ✓ Aucune blocage critique signalé
    \item ✓ Documentation complète et satisfaisante
\end{itemize}

% =========== CHAPITRE 10: SIGNATURE ===========
\chapter*{Approbations}

\vspace{2cm}

\begin{table}[H]
\centering
\begin{tabularx}{\textwidth}{|l|l|l|}
\hline
\rowcolor{LightGray}
\textbf{Rôle} & \textbf{Signature} & \textbf{Date} \\
\hline
Chef de Projet & & \\
\hline
 & & \\
\hline
Sponsor Exécutif & & \\
\hline
 & & \\
\hline
Architecte Technique & & \\
\hline
 & & \\
\hline
Responsable QA & & \\
\hline
 & & \\
\hline
\end{tabularx}
\end{table}

\vspace{1cm}

\noindent Date d'approbation finale : \underline{\hspace{4cm}}

% =========== ANNEXES ===========
\appendix

\chapter{Glossaire}

\begin{description}
    \item[API] Application Programming Interface
    \item[JWT] JSON Web Token
    \item[REST] Representational State Transfer
    \item[RBAC] Role-Based Access Control
    \item[QA] Quality Assurance
    \item[SLA] Service Level Agreement
    \item[DevOps] Development and Operations
    \item[CI/CD] Continuous Integration/Continuous Deployment
    \item[ROI] Return on Investment
    \item[UI/UX] User Interface/User Experience
\end{description}

\chapter{Références et Documentation}

\section{Documentation Technique}

\begin{itemize}
    \item Spring Boot Official Documentation : \url{https://spring.io/projects/spring-boot}
    \item React Documentation : \url{https://react.dev}
    \item Material-UI Documentation : \url{https://mui.com}
    \item PostgreSQL Documentation : \url{https://www.postgresql.org/docs}
    \item Docker Documentation : \url{https://docs.docker.com}
\end{itemize}

\section{Fichiers de Projet}

\begin{itemize}
    \item Repository GitHub : À compléter
    \item Figma Design : \url{https://www.figma.com/design/8OrQuigMzGUMzQadYRmksk/SolarEase-Web-App-Design}
    \item Jira/Trello Board : À compléter
    \item Wiki/Confluence : À compléter
\end{itemize}

\chapter{API REST Endpoints}

\section{Gateway Service}

\begin{lstlisting}[language=bash]
BASE_URL: http://localhost:8080/api

Authentication:
  POST   /auth/login          - Connexion utilisateur
  POST   /auth/register       - Inscription nouvel utilisateur
  POST   /auth/refresh        - Rafraîchir token JWT
  POST   /auth/verify-otp     - Vérifier code OTP

Identity Service:
  GET    /users/{id}          - Récupérer profil utilisateur
  PUT    /users/{id}          - Mettre à jour profil
  GET    /users               - Lister tous les utilisateurs
  POST   /roles               - Créer nouveau rôle

Dimensioning Service:
  POST   /dimensioning/calculate        - Calculer dimensionnement
  GET    /dimensioning/project/{id}     - Récupérer dimensionnements d'un projet
  GET    /dimensioning/{id}             - Récupérer dimensionnement spécifique

Project Service:
  POST   /projects            - Créer nouveau projet
  GET    /projects            - Lister tous les projets
  GET    /projects/{id}       - Détails d'un projet
  PUT    /projects/{id}       - Mettre à jour projet
  GET    /clients             - Lister les clients
  POST   /clients             - Ajouter nouveau client
  GET    /equipment           - Lister les équipements
\end{lstlisting}

\chapter{Exigences Non-Fonctionnelles}

\section{Performance}

\begin{itemize}
    \item Temps de chargement initial de l'app < 3 secondes
    \item Temps de réponse API < 500ms (95e percentile)
    \item Support concurrent : 1000+ utilisateurs simultanés
    \item Throughput : 500+ requests/seconde
\end{itemize}

\section{Disponibilité}

\begin{itemize}
    \item SLA : 99.5\% uptime
    \item RTO (Recovery Time Objective) : < 1 heure
    \item RPO (Recovery Point Objective) : < 15 minutes
    \item Backup automatique : Quotidien
\end{itemize}

\section{Maintenabilité}

\begin{itemize}
    \item Code modulaire et découplé
    \item Convention de nommage consistent
    \item Documentation en ligne
    \item Tests automatisés
\end{itemize}

\section{Évolutivité}

\begin{itemize}
    \item Architecture microservices horizontalement scalable
    \item Base de données optimisée pour croissance
    \item Caching stratégique (Redis si nécessaire)
    \item Load balancing en place
\end{itemize}

\section{Intégration de la Gestion des Ressources d'Architecture (RAG)}

L'intégration de la gestion des ressources d'architecture (RAG) dans le service de dimensionnement permet d'optimiser l'utilisation des ressources tout au long du processus de dimensionnement solaire. Cela inclut :

\begin{itemize}
    \item \textbf{Planification des Ressources} : Évaluation des besoins en ressources pour chaque projet de dimensionnement.
    \item \textbf{Suivi des Ressources} : Outils pour suivre l'utilisation des ressources en temps réel.
    \item \textbf{Optimisation des Ressources} : Stratégies pour maximiser l'efficacité des ressources allouées.
    \item \textbf{Rapports et Analyses} : Génération de rapports réguliers sur l'état des ressources et leur utilisation.
\end{itemize}

Cette intégration vise à garantir que les projets de dimensionnement sont réalisés de manière efficace et avec un minimum de gaspillage de ressources.

\end{document}
